% ======================================================
% 第 41 天
% ======================================================
\setday{41}{一阶线性微分方程与伯努利方程}

\oneproblem{
    试写出以 $y=x^2-e^x$ 和 $y=x^2$ 为特解的一阶非齐次线性微分方程.
}

\oneproblem{
    已知可导函数 $f(x)$ 满足
    $
    f(x)\cos x + 2\int_0^x f(t)\sin t\,\mathrm{d}t = x+1,
    $
    试求 $f(x)$ 的表达式.
}

\oneproblem{
    设 $y=e^x$ 是微分方程 $xy'+p(x)y=x$ 的一个解，求此微分方程满足条件 $y|_{x=\ln2}=0$ 的特解.
}

\oneproblem{
    设曲线 $L$ 位于 $xOy$ 面的第一象限内，$L$ 上任一点 $M$ 处的切线与 $y$ 轴总相交，交点记 $A$，已知 $|\overline{MA}|=|\overline{OA}|$，且 $L$ 过点 $\left(\dfrac{3}{2},\dfrac{3}{2}\right)$，求 $L$ 的方程.
}

\oneproblem{
    设函数 $f(x)$ 在闭区间 $[0,1]$ 上连续，在开区间 $(0,1)$ 内大于零，并满足
    $
    xf'(x)=f(x)+\dfrac{3a}{2}x^2\ (a\ \text{为常数}),
    $
    又曲线 $y=f(x)$ 与 $x=1,y=0$ 所围成的图形 $S$ 的面积值为 $2$，求函数 $y=f(x)$，并问 $a$ 为何值时，图形 $S$ 绕 $x$ 轴旋转一周所得的旋转体的体积最小.
}

\oneproblem{
    设 $y(x)$ 是微分方程 $2xy'-4y=2\ln x-1$ 满足 $y(1)=\dfrac{1}{4}$ 的解，求曲线 $y=y(x)\ (1\le x\le e)$ 的弧长.
}

\oneproblem{
    设 $f(x)$ 为连续函数.
    \begin{enumerate}[label=(\arabic*)]
        \item 求初值问题 $\begin{cases} y'+ay=f(x), \\ y|_{x=0}=0 \end{cases}$ 的解 $y=y(x)$，其中 $a$ 是正常数；
        \item 若 $|f(x)|\le k$（$k$ 为常数），证明：当 $x\ge 0$ 时，有 $|y(x)|\le \dfrac{k}{a}(1-e^{-ax})$.
    \end{enumerate}
}

\oneproblem{
    设有微分方程 $y'-2y=\varphi(x)$，其中 $\varphi(x)=\begin{cases} 2 & x<1, \\ 0 & x>1 \end{cases}$，试求在 $(-\infty,+\infty)$ 内的连续函数 $y=y(x)$，使之在 $(-\infty,1)$ 和 $(1,+\infty)$ 内都满足所给方程，且 $y(0)=0$.
}

\oneproblem{
    求微分方程 $x\mathrm{d}y+(x-2y)\mathrm{d}x=0$ 的一个解 $y=y(x)$，使得由曲线 $y=y(x)$ 与直线 $x=1,x=2$ 以及 $x$ 轴所围成的平面图形绕 $x$ 轴旋转一周的旋转体体积最小.
}

\oneproblem{
    设 $y=f(x)$ 是第一象限内连接点 $A(0,1),B(1,0)$ 的一段连续曲线，$M(x,y)$ 为该曲线上任意一点，点 $C$ 为 $M$ 在 $x$ 轴上的投影，$O$ 为坐标原点。若梯形 $OCMA$ 的面积与曲边三角形 $CBM$ 的面积之和为 $\dfrac{x^3}{6}+\dfrac{1}{3}$，求 $f(x)$ 的表达式.
}

\oneproblem{
    设 $F(x)=f(x)g(x)$，其中函数 $f(x),g(x)$ 在 $(-\infty,+\infty)$ 内满足以下条件：
    $
    f'(x)=g(x),\ g'(x)=f(x),\ \text{且}\ f(0)=0,\ f(x)+g(x)=2e^x.
    $
    \begin{enumerate}[label=(\arabic*)]
        \item 求 $F(x)$ 所满足的一阶微分方程；
        \item 求出 $F(x)$ 的表达式.
    \end{enumerate}
}

\oneproblem{
    设曲线 $y=f(x)$，其中 $f(x)$ 是可导函数，且 $f(x)>0$。已知曲线 $y=f(x)$ 与直线 $y=0,x=1$ 及 $x=t\ (t>1)$ 所围成的曲边梯形绕 $x$ 轴旋转一周所得的立体体积值是该曲边梯形面积值的 $\pi t$ 倍，求该曲线的方程.
}

\oneproblem{
    已知连续函数 $f(x)$ 满足
    $
    \int_0^x f(t)\mathrm{d}t + \int_0^x tf(x-t)\mathrm{d}t = ax^2.
    $
    \begin{enumerate}[label=(\arabic*)]
        \item 求 $f(x)$；
        \item 若 $f(x)$ 在区间 $[0,1]$ 上的平均值为 $1$，求 $a$ 的值.
    \end{enumerate}
}

\oneproblem{
    已知微分方程 $y'+y=f(x)$，其中 $f(x)$ 是 $\mathbb{R}$ 上的连续函数.
    \begin{enumerate}[label=(\arabic*)]
        \item 当 $f(x)=x$ 时，求微分方程的通解；
        \item 当 $f(x)$ 为周期函数时，证明微分方程有通解与其对应，且该通解也为周期函数.
    \end{enumerate}
}

\oneproblem{
    已知 $y(x)$ 满足微分方程
    $
    y'-xy=\dfrac{1}{2\sqrt{x}}e^{\frac{x^2}{2}},
    $
    且有 $y(1)=\sqrt{e}$.
    \begin{enumerate}[label=(\arabic*)]
        \item 求 $y(x)$；
        \item 设平面区域 $D=\{(x,y)\mid 1\le x\le 2,0\le y\le y(x)\}$，求平面区域 $D$ 绕 $x$ 轴旋转所成旋转体体积.
    \end{enumerate}
}

\oneproblem{
    函数 $y=y(x)\ (x>0)$ 满足微分方程 $xy'-6y=-6$，且 $y(\sqrt{3})=10$.
    \begin{enumerate}[label=(\arabic*)]
        \item 求 $y(x)$；
        \item $P$ 为曲线 $y=y(x)$ 上一点，曲线 $y=y(x)$ 在点 $P$ 的法线在 $y$ 轴上的截距为 $I_y$，为使得 $I_y$ 最小，求 $P$ 的坐标.
    \end{enumerate}
}

\oneproblem{
    设 $y=y(x)$ 满足 $y'+\dfrac{1}{2\sqrt{x}}y=2+\sqrt{x},y(1)=3$，求曲线 $y=y(x)$ 的渐近线.
}

\oneproblem{
    设 $y(x)$ 是微分方程 $2xy'-4y=2\ln x-1$ 满足 $y(1)=\dfrac{1}{4}$ 的解，求曲线 $y=y(x)\ (1\le x\le e)$ 的弧长.
}

\oneproblem{
    设有微分方程 $y'+y=f(x)$，其中
    $
    f(x)=\begin{cases} 2, & 0\le x\le 1, \\ 0, & x>1. \end{cases}
    $
    试求此方程满足初始条件 $y|_{x=0}=0$ 的连续解.
}

\oneproblem{
    设函数 $y(x)$ 满足微分方程 $\dfrac{\mathrm{d}y}{\mathrm{d}x}=\dfrac{xy^2+\sin x}{2y}$。试求函数 $g(x)=e^{-\frac{x^2}{2}}y^2(x)$ 在 $[-\pi,\pi]$ 内极值点，并判定极值点的类型.
}

\oneproblem{
    \begin{enumerate}[label=(\arabic*)]
        \item 求解微分方程 $\begin{cases} \dfrac{\mathrm{d}y}{\mathrm{d}x}-xy=xe^{x^2}, \\ y(0)=1. \end{cases}$
        \item 如 $y=f(x)$ 为上述方程的解，证明：$\displaystyle \lim_{n\to\infty}\int_0^1 \dfrac{nf(x)}{n^2x^2+1}\,\mathrm{d}x=\dfrac{\pi}{2}$.
    \end{enumerate}
}

% ======================================================
% 第 42 天
% ======================================================
\setday{42}{可降阶的高阶微分方程}

\oneproblem{
    求下列微分方程的通解.
    \begin{enumerate}[label=(\arabic*)]
        \item $y''-(y')^3=0$.
        \item $xyy''+x(y')^2=3yy'$.
        \item $y'\cdot y'''=3(y'')^2$.
    \end{enumerate}
}

\oneproblem{
    求解下列初值问题：
    \begin{enumerate}[label=(\arabic*)]
        \item $\begin{cases} y''-a(y')^2=0, \\ y(0)=0,y'(0)=-1 \end{cases}\ (a\ne 0)$.
        \item $\begin{cases} y''(x+y'^2)=y', \\ y(1)=y'(1)=1. \end{cases}$
    \end{enumerate}
}

\oneproblem{
    设对任意 $x>0$，曲线 $y=f(x)$ 上点 $(x,f(x))$ 处的切线在 $y$ 轴上的截距等于 $\dfrac{1}{x}\int_0^x f(t)\,\mathrm{d}t$，求 $f(x)$ 的一般表达式.
}

\oneproblem{
    已知函数 $f(x)$ 可导，且 $f'(x)>0\ (x>0)$。曲线 $y=f(x)\ (x\ge 0)$ 经过坐标原点 $O$，其上任意一点 $M$ 处的切线与 $x$ 轴相交于点 $T$，过点 $M$ 做 $MP$ 垂直于 $x$ 轴于点 $P$，且曲线 $y=f(x)$、直线 $MP$ 以及 $x$ 轴所围成图形的面积与三角形 $MTP$ 的面积比恒为 $3:2$，求满足上述条件的曲线的方程.
}

\oneproblem{
    在上半平面求一条向上凹的曲线，其上任一点 $P(x,y)$ 处的曲率等于此曲线在该点的法线段 $PQ$ 长度的倒数（$Q$ 是法线与 $x$ 轴的交点），且曲线在点 $(1,1)$ 处的切线与 $x$ 轴平行.
}

\oneproblem{
    设函数 $y(x)\ (x\ge 0)$ 二阶可导且 $y'(x)>0,y(0)=1$。过曲线 $y=f(x)$ 上任意一点 $P(x,y)$ 作该曲线的切线及 $x$ 轴的垂线，上述两直线与 $x$ 轴所围成的三角形的面积记为 $S_1$，区间 $[0,x]$ 上以 $y=y(x)$ 为曲边的曲边梯形面积记为 $S_2$，并设 $2S_1-S_2$ 恒为 $1$，求此曲线 $y=y(x)$ 的方程.
}

\oneproblem{
    设函数 $y(x)$ 具有二阶导数，且曲线 $l:y=y(x)$ 与直线 $y=x$ 相切于原点，记 $\alpha$ 为曲线 $l$ 在点 $(x,y)$ 处切线的倾角，若 $\dfrac{\mathrm{d}\alpha}{\mathrm{d}x}=\dfrac{\mathrm{d}y}{\mathrm{d}x}$，求 $y(x)$ 的表达式.
}

\oneproblem{
    设 $y=f(x)$ 由参数方程 $\begin{cases} x=2t+t^2 \\ y=\psi(t) \end{cases}\ (t>-1)$ 所确定。且 $\dfrac{\mathrm{d}^2y}{\mathrm{d}x^2}=\dfrac{3}{4(1+t)}$，其中 $\psi(t)$ 具有二阶导数，曲线 $y=\psi(t)$ 与 $y=\int_1^{t^2}e^{-u^2}\mathrm{d}u+\dfrac{3}{2e}$ 在 $t=1$ 处相切. 求函数 $\psi(t)$.
}

\oneproblem{
    函数 $f(x)$ 在 $[0,+\infty)$ 上可导，$f(0)=1$，且满足等式
    $
    f'(x)+f(x)-\dfrac{1}{x+1}\int_0^x f(t)\,\mathrm{d}t=0.
    $
    \begin{enumerate}[label=(\arabic*)]
        \item 求导数 $f'(x)$；
        \item 证明：当 $x\ge 0$ 时，不等式 $e^{-x}\le f(x)\le 1$ 成立.
    \end{enumerate}
}

% ======================================================
% 第 43 天
% ======================================================
\setday{43}{线性微分方程解的结构与二阶变系数线性微分方程求解}

\oneproblem{
    设 $y_1,y_2$ 是一阶线性非齐次微分方程 $y'+p(x)y=q(x)$ 的两个特解，若常数 $\lambda,\mu$ 使 $\lambda y_1+\mu y_2$ 是该方程的解，$\lambda y_1-\mu y_2$ 是该方程对应的齐次方程的解，则（ \quad ）
    \begin{multicols}{2}
    \begin{enumerate}[label=(\Alph*)]
        \item $\lambda=\dfrac{1}{2},\mu=\dfrac{1}{2}$
        \item $\lambda=-\dfrac{1}{2},\mu=-\dfrac{1}{2}$
        \item $\lambda=\dfrac{2}{3},\mu=\dfrac{1}{3}$
        \item $\lambda=\dfrac{2}{3},\mu=\dfrac{2}{3}$
    \end{enumerate}
    \end{multicols}
}

\oneproblem{
    设线性无关的函数 $y_1,y_2,y_3$ 都是二阶非齐次线性微分方程 $y''+p(x)y'+q(x)y=f(x)$ 的解，$C_1,C_2$ 均为任意常数，则该非齐次微分方程的通解是（ \quad ）
    \begin{multicols}{2}
    \begin{enumerate}[label=(\Alph*)]
        \item $C_1y_1+C_2y_2+y_3$
        \item $C_1y_1+C_2y_2-(1-C_1+C_2)y_3$
        \item $C_1y_1+C_2y_2-(1-C_1-C_2)y_3$
        \item $C_1y_1+C_2y_2+(1-C_1-C_2)y_3$
    \end{enumerate}
    \end{multicols}
}

\oneproblem{
    已知 $y_1=xe^x+e^{2x},y_2=xe^x+e^{-x},y_3=xe^x+e^{2x}+e^{-x}$ 是某二阶线性非齐次微分方程的三个解，求此微分方程.
}

\oneproblem{
    设非齐次线性微分方程 $y'+P(x)y=Q(x)$ 有两个不同的解 $y_1(x),y_2(x),C$ 为任意常数，则该方程的通解是（ \quad ）
    \begin{multicols}{2}
    \begin{enumerate}[label=(\Alph*)]
        \item $C[y_1(x)-y_2(x)]$
        \item $y_1(x)+C[y_1(x)-y_2(x)]$
        \item $C[y_1(x)+y_2(x)]$
        \item $y_1(x)+C[y_1(x)+y_2(x)]$
    \end{enumerate}
    \end{multicols}
}

\oneproblem{
    已知 $y_1(x)=e^x,y_2(x)=u(x)e^x$ 是二阶微分方程 $(2x-1)y''-(2x+1)y'+2y=0$ 的解，若 $u(-1)=e,u(0)=-1$，求 $u(x)$，并写出该微分方程的通解.
}

\oneproblem{
    设微分方程 $y''+\dfrac{1}{x}y'-q(x)y=0$ 有两个特解 $y_1(x),y_2(x)$，且 $y_1\cdot y_2=1$。求此方程中的 $q(x)$，并求该微分方程的通解.
}

\oneproblem{
    求微分方程 $y''-\dfrac{x+2}{x}y'+\dfrac{x+2}{x^2}y=x^2$ 的通解.
}

\oneproblem{
    解微分方程 $(3t^3+t)\dfrac{\mathrm{d}^2x}{\mathrm{d}t^2}+2\dfrac{\mathrm{d}x}{\mathrm{d}t}-6tx=4-12t^2$，已知方程有两个解：$x_1(t)=2t,x_2(t)=(1+t)^2$.
}

\oneproblem{
    设函数 $y=f(x)$ 满足微分方程 $y''-3y'+2y=2e^x$，且其图形在点 $(0,1)$ 处的切线与曲线 $y=x^2-x+1$ 在该点的切线重合，求函数 $f(x)$.
}

\oneproblem{
    设二阶线性方程 $\dfrac{\mathrm{d}^2y}{\mathrm{d}x^2}+P(x)\dfrac{\mathrm{d}y}{\mathrm{d}x}+Q(x)y=f(x)$，其中 $P(x)$ 和 $Q(x)$ 都是某区间上的连续函数.
    \begin{enumerate}[label=(\arabic*)]
        \item 证明：存在变换 $y=u(x)v(x)$，可以把给定方程化为方程 $\dfrac{\mathrm{d}^2u}{\mathrm{d}x^2}+R(x)u=f(x)e^{\int\frac{1}{2}P(x)\mathrm{d}x}$，并求出 $R(x)$.
        \item 求微分方程 $xy''+2y'-xy=e^x$ 的通解.
    \end{enumerate}
}

\oneproblem{
    已知函数 $\sin^2x,\cos^2x$ 都是方程 $y''+P(x)y'+Q(x)y=0$ 的解，
    \begin{enumerate}[label=(\arabic*)]
        \item 证明：$\sin^2x,\cos^2x$ 构成基本解组；
        \item 证明：$1,\cos2x$ 也构成基本解组；
        \item 求 $P(x),Q(x)$；
        \item 求 $y''+P(x)y'+Q(x)y=\sin2x$ 的通解.
    \end{enumerate}
}

% ======================================================
% 第 44 天
% ======================================================
\setday{44}{常系数齐次线性微分方程}

\oneproblem{
    具有特解 $y_1=e^{-x}, y_2=2xe^{-x}, y_3=3e^x$ 的 3 阶常系数齐次线性微分方程是（ \quad ）
    \begin{multicols}{2}
    \begin{enumerate}[label=(\Alph*)]
        \item $y'''-y''-y'+y=0$
        \item $y'''+y''-y'-y=0$
        \item $y'''-6y''+11y'-6y=0$
        \item $y'''-2y''-y'+2y=0$
    \end{enumerate}
    \end{multicols}
}

\oneproblem{
    设 $y=e^x(C_1\sin x+C_2\cos x)$（$C_1,C_2$ 为任意常数）为某二阶常系数线性齐次微分方程的通解，试求该微分方程.
}

\oneproblem{
    已知 $y_1=e^x$ 和 $y_2=xe^x$ 是齐次二阶常系数线性微分方程的解，试求该微分方程.
}

\oneproblem{
    试求三阶常系数线性齐次微分方程 $y'''-2y''+y'-2y=0$ 的通解.
}

\oneproblem{
    设函数 $f(x)$ 满足在
    $
    \begin{cases}
    f''(x)+af'(x)+f(x)=0\ (a>0), \\
    f(0)=m,\ f'(0)=n,
    \end{cases}
    $
    求广义积分 $\displaystyle \int_0^{+\infty} f(x)\,\mathrm{d}x$.
}

\oneproblem{
    设 $y=y(x)$ 满足条件
    $
    \begin{cases}
    y''+4y'+4y=0, \\
    y(0)=2,\ y'(0)=-4,
    \end{cases}
    $
    求广义积分 $\displaystyle \int_0^{+\infty} y(x)\,\mathrm{d}x$.
}

% ======================================================
% 第 45 天
% ======================================================
\setday{45}{常系数非齐次线性微分方程}

\oneproblem{
    微分方程 $y''-y=e^x+1$ 的一个特解应具有形式为（以下 $a,b$ 为常数）（ \quad ）
    \begin{multicols}{2}
    \begin{enumerate}[label=(\Alph*)]
        \item $ae^x+b$
        \item $axe^x+b$
        \item $ae^x+bx$
        \item $axe^x+bx$
    \end{enumerate}
    \end{multicols}
}

\oneproblem{
    设 $y=y(x)$ 是二阶常系数微分方程 $y''+py'+qy=e^{3x}$ 满足初始条件 $y(0)=y'(0)=0$ 的特解，则当 $x\to0$ 时，函数 $\dfrac{\ln(1+x^2)}{y(x)}$ 的极限（ \quad ）
    \begin{multicols}{2}
    \begin{enumerate}[label=(\Alph*)]
        \item 不存在
        \item 等于 1
        \item 等于 2
        \item 等于 3
    \end{enumerate}
    \end{multicols}
}

\oneproblem{
    微分方程 $y''+y=x^2+1+\sin x$ 的特解形式可设为（ \quad ）
    \begin{multicols}{2}
    \begin{enumerate}[label=(\Alph*)]
        \item $y^*=ax^2+bx+c+x(A\sin x+B\cos x)$
        \item $y^*=x(ax^2+bx+c+A\sin x+B\cos x)$
        \item $y^*=ax^2+bx+c+A\sin x$
        \item $y^*=ax^2+bx+c+B\cos x$
    \end{enumerate}
    \end{multicols}
}

\oneproblem{
    求下列微分方程的通解：
    \begin{enumerate}[label=(\arabic*)]
        \item $y^{(4)}+y''=x+e^x+3\sin x$.
        \item $y''-y'=x\cdot2^x+\cos x$.
        \item $y''-y=\cos^3x$.
        \item $y''+4y'+4y=e^{ax}$，其中 $a$ 为实数.
    \end{enumerate}
}

\oneproblem{
    用变量代换 $x=\cos t\ (0<t<\pi)$ 化简微分方程 $(1-x^2)y''-xy'+y=0$，并求其满足 $y|_{x=0}=1,\ y'|_{x=0}=2$ 的特解.
}

\oneproblem{
    设函数 $y(x)$ 满足方程 $y''+2y'+ky=0$，其中 $0<k<1$.
    \begin{enumerate}[label=(\arabic*)]
        \item 证明：反常积分 $\displaystyle \int_0^{+\infty} y(x)\,\mathrm{d}x$ 收敛；
        \item 若 $y(0)=1,\ y'(0)=1$，求 $\displaystyle \int_0^{+\infty} y(x)\,\mathrm{d}x$ 的值.
    \end{enumerate}
}

\oneproblem{
    设函数 $y=f(x)$ 满足 $y''+2y'+5y=0$，且有 $f(0)=1,\ f'(0)=-1$.
    \begin{enumerate}[label=(\arabic*)]
        \item 求 $f(x)$ 的表达式；
        \item 设 $a_n=\displaystyle \int_{n\pi}^{+\infty} f(x)\,\mathrm{d}x$，求 $\displaystyle \sum_{n=1}^\infty a_n$.
    \end{enumerate}
}

\oneproblem{
    已知函数 $f(x)$ 满足方程 $f''(x)+f'(x)-2f(x)=0$ 及 $f''(x)+f(x)=2e^x$.
    \begin{enumerate}[label=(\arabic*)]
        \item 求 $f(x)$ 的表达式；
        \item 求曲线 $y=f(x^2)\displaystyle \int_0^x f(-t^2)\,\mathrm{d}t$ 的拐点.
    \end{enumerate}
}

\oneproblem{
    已知方程 $y''-y'-2y=-2x-1$ 满足 $\begin{cases} y(0)=0 \\ y'(0)=0 \end{cases}$ 与 $\begin{cases} y(0)=1 \\ y'(0)=2 \end{cases}$ 的两个特解恰好是微分方程 $(x^2+1)y'''-2xy'-4(x^2-x+1)y=-2(2x^3-2x^2+3x)$ 的两个特解，求该微分方程的通解.
}

% ======================================================
% 第 46 天
% ======================================================
\setday{46}{特殊微分方程解法举例}

\oneproblem{
    求下列微分方程的通解：
    \begin{enumerate}[label=(\arabic*)]
        \item $(x+2)^2y''-2(x+2)y'+2y=x$.
        \item $(y')^2+xy'-y=0$.
        \item $(2x-1)^2\dfrac{\mathrm{d}^2y}{\mathrm{d}x^2}-4(2x-1)\dfrac{\mathrm{d}y}{\mathrm{d}x}+8y=8x$.
        \item $y''\sin x+2y'\cos x+3y\sin x=e^x$.
    \end{enumerate}
}

\oneproblem{
    求在第一象限中的一条曲线，使其上每一点处的切线与两坐标轴所围成的三角形的面积均等于 2.
}

\oneproblem{
    求微分方程 $y''+y'-2y=\dfrac{e^x}{1+e^x}$ 的通解.
}

\oneproblem{
    设 $f(x)$ 在 $[0,+\infty)$ 上是有界连续函数，证明：方程 $y''+14y'+13y=f(x)$ 的每一个解在 $[0,+\infty)$ 上都是有界函数.
}

\oneproblem{
    用参数方程法求下列微分方程的通解：
    \begin{enumerate}[label=(\arabic*)]
        \item $x^3+y'^3-3xy'=0$.
        \item $(y')^3-y^2(a-y')=0$.
        \item $y=\ln\sqrt{1+y'^2}$.
        \item $(x^2+y^2)y''=2(1+y'^2)(xy'-y)$.
    \end{enumerate}
}

\oneproblem{
    设函数 $y=y(x)$ 满足
    $
    xy+\int_1^x\left[3y+t^2y''(t)\right]\mathrm{d}t=5\ln x,\ x\ge1\ \text{且}\ y'(1)=0,
    $
    求函数 $y(x)$ 的表达式.
}

\oneproblem{
    \begin{enumerate}[label=(\arabic*)]
        \item 试将 $x=x(y)$ 所满足的微分方程 $\dfrac{\mathrm{d}^2x}{\mathrm{d}y^2}+(y+\sin x)\left(\dfrac{\mathrm{d}x}{\mathrm{d}y}\right)^3=0$ 变换为 $y=y(x)$ 满足的微分方程.
        \item 求变换后的微分方程满足初始条件 $y(0)=0,\ y'(0)=\dfrac{3}{2}$ 的解.
    \end{enumerate}
}

\oneproblem{
    设函数 $y=y(x)$ 是区间 $(-\pi,\pi)$ 内过点 $\left(-\dfrac{\pi}{\sqrt{2}},\dfrac{\pi}{\sqrt{2}}\right)$ 的光滑曲线，当 $-\pi<x<0$ 时，曲线上任一点处的法线都过原点，当 $0\le x<\pi$ 时，函数 $y(x)$ 满足 $y''+y+x=0$. 求 $y(x)$ 的表达式.
}

\oneproblem{
    求 $y''+4y=3|\sin x|$ 在 $[-\pi,\pi]$ 上满足 $y\left(\dfrac{\pi}{2}\right)=0,\ y'\left(\dfrac{\pi}{2}\right)=1$ 的特解.
}

% ======================================================
% 第 47 天
% ======================================================
\setday{47}{微分方程的建模及其应用}

\oneproblem{
    有高为 $H$ cm 的半球容器，水从它的底部小孔流出，小孔的横截面积为 $S$ cm$^2$，开始时容器盛满水，问多久水会流完？已知水高度为 $h$ 的小孔水流出的速度为 $0.62\sqrt{2gh}$ (cm/s).
}

\oneproblem{
    一个半球体状的雪堆，其体积融化的速率与半球面面积 $S$ 成正比，比例常数 $K>0$。假设在融化过程中雪堆始终保持半球体状，已知半径为 $r_0$ 的雪堆在开始融化的 3 小时内，融化了其体积的 $\dfrac{7}{8}$，问雪堆全部融化需要多少小时？
}

\oneproblem{
    设单位质点在水平面内作直线运动，初速度 $v|_{t=0}=v_0$，已知阻力与速度成正比（比例常数为 1），问 $t$ 为多少时此质点的速度为 $\dfrac{v_0}{3}$？并求此时刻该质点所经过的路程.
}

\oneproblem{
    在某一人群中推广新技术是通过其中已掌握新技术的人进行的。设该人群的总人数为 $N$，在 $t=0$ 时刻已掌握新技术人数为 $x_0$，在任意时刻 $t$ 已掌握新技术的人数为 $x(t)$（将 $x(t)$ 视为连续可微变量），其变化率与已掌握新技术人数和未掌握新技术人数之积成正比，比例常数 $k>0$，求 $x(t)$.
}

\oneproblem{
    从船上向海中沉放某种探测仪器，按探测要求，需确定仪器的下沉深度 $y$（从海平面算起）与下沉速度 $v$ 之间的函数关系。设仪器在重力作用下，从海平面由静止开始铅直下沉，在下沉过程中还受到阻力和浮力的作用。设仪器的质量为 $m$，体积为 $B$，海水比重为 $\rho$，仪器所受的阻力与下沉速度成正比，比例系数为 $k(k>0)$。试建立 $y$ 与 $v$ 所满足的微分方程，并求出函数关系式 $y=y(v)$.
}

\oneproblem{
    某种飞机在机场降落时，为了减少滑行距离，在触地的瞬间，飞机尾部张开减速伞以增大阻力，使飞机迅速减速并停下。现有一质量为 9000 kg 的飞机，着陆时的水平速度为 700 km/h。经测试，减速伞打开后，飞机所受的总阻力与飞机的速度成正比（比例系数为 $k=6.0\times10^6$）问从着陆点算起，飞机滑行的最长距离是多少？（kg 表示千克，km/h 表示千米/小时）.
}

\oneproblem{
    有一平底容器，其内侧壁是由曲线 $x=\varphi(y)\ (y\ge0)$ 绕 $y$ 轴旋转而成的旋转曲面，容器的底面圆的半径为 2m。根据设计要求，当以 $3\,\text{m}^3/\text{min}$ 的速率向容器内注入液体时，液面的面积将以 $\pi\,\text{m}^2/\text{min}$ 的速率均匀扩大（假设注入液体前，容器内无液体）.
    \begin{enumerate}[label=(\arabic*)]
        \item 根据 $t$ 时刻液面的面积，写出 $t$ 与 $\varphi(y)$ 之间的关系式；
        \item 求曲线 $x=\varphi(y)$ 的方程.（注：m 表示长度单位米，min 表示时间单位分.）
    \end{enumerate}
}

\oneproblem{
    某湖泊的水量为 $V$，每年排入湖泊内含污染物 $A$ 的污水量为 $\dfrac{V}{6}$，流入湖泊内不含 $A$ 的水量为 $\dfrac{V}{6}$，流出湖泊的水量为 $\dfrac{V}{3}$。已知 1999 年年底湖中 $A$ 的含量为 $5m_0$，超过国家规定指标，为了治理污水，从 2000 年年初起，限定排入湖泊中含 $A$ 污水的浓度不超过 $\dfrac{m_0}{V}$。问至多需经过多少年，湖泊中污染物 $A$ 的含量将至 $m_0$ 以内？（注：设湖水中 $A$ 的浓度是均匀的.）
}

\oneproblem{
    已知高温物体置于低温介质中，任一时刻物体温度对时间的变化率与该时刻物体和介质的温差成正比，现将一初始温度为 $120^\circ\text{C}$ 的物体在 $20^\circ\text{C}$ 的恒温介质中冷却，30 min 后该物体温度降至 $30^\circ\text{C}$，若要使物体的温度继续降至 $21^\circ\text{C}$，还需冷却多长时间？
}

\oneproblem{
    质量为 $m$ 的质点受力 $F$ 的作用沿 $Ox$ 轴作直线运动。设力 $F$ 仅是时间 $t$ 的函数：$F=F(t)$。在开始时刻 $t=0$ 时，$F(0)=F_0$，随着时间 $t$ 的增大，此力 $F$ 均匀地减小，直到 $t=T$ 时 $F(T)=0$。如果开始时质点在原点且初速度为 0，求这质点的运动规律.
}

\oneproblem{
    设有一均匀、柔软的绳索，两端固定，绳索仅受到重力的作用下垂，试问该绳索在平衡状态下是怎样的曲线？
}

\oneproblem{
    设物体 $A$ 从点 $(0,1)$ 出发，以大小为常数 $v$ 的速度沿 $y$ 轴正向运动，同一时刻，物体 $B$ 从 $(-1,0)$ 出发，速度大小为 $2v$，方向指向 $A$，试求物体 $B$ 在 $y$ 轴左侧的运动轨迹方程.
}

\oneproblem{
    设某一海湾的海岸线呈直角形，设直角顶点为原点，海岸边界线为两坐标轴，则海水的流线为双曲线族 $xy=C$（$C$ 为常数）。有一条船在海湾内行驶，船身保持和水流方向成 30 度角并指向海岸，试给出该船行驶的路线方程.
}

% ======================================================
% 第 48 天
% ======================================================
\setday{48}{向量代数：向量及其运算}

\oneproblem{
    设 $\vec{A},\vec{B},\vec{C}$ 为非零矢量，且 $\vec{A}\cdot\vec{B}=0,\ \vec{A}\times\vec{C}=0$，则（ \quad ）
    \begin{multicols}{2}
    \begin{enumerate}[label=(\Alph*)]
        \item $\vec{A}\parallel\vec{B}$ 且 $\vec{B}\perp\vec{C}$
        \item $\vec{A}\perp\vec{B}$ 且 $\vec{B}\parallel\vec{C}$
        \item $\vec{A}\parallel\vec{C}$ 且 $\vec{B}\perp\vec{C}$
        \item $\vec{A}\perp\vec{C}$ 且 $\vec{B}\parallel\vec{C}$
    \end{enumerate}
    \end{multicols}
}

\oneproblem{
    若非零向量 $\mathbf{a},\mathbf{b}$ 满足关系式 $|\mathbf{a}-\mathbf{b}|=|\mathbf{a}+\mathbf{b}|$，则必有（ \quad ）
    \begin{multicols}{2}
    \begin{enumerate}[label=(\Alph*)]
        \item $\mathbf{a}-\mathbf{b}=\mathbf{a}+\mathbf{b}$
        \item $\mathbf{a}=\mathbf{b}$
        \item $\mathbf{a}\cdot\mathbf{b}=0$
        \item $\mathbf{a}\times\mathbf{b}=\mathbf{0}$
    \end{enumerate}
    \end{multicols}
}

\oneproblem{
    设 $(\vec{a}\times\vec{b})\cdot\vec{c}=2$，则 $[(\vec{a}+\vec{b})\times(\vec{b}+\vec{c})]\cdot(\vec{c}+\vec{a})=$ \underline{\hspace{2cm}}.
}

\oneproblem{
    设点 $A$ 位于第一卦限，向径 $\overrightarrow{OA}$ 与 $x$ 轴、$y$ 轴的夹角依次为 $\dfrac{\pi}{3},\dfrac{\pi}{4}$，且 $|OA|=6$，求点 $A$ 的坐标.
}

\oneproblem{
    已知向量 $\vec{a}=(2,-3,6),\vec{b}=(-1,2,-2)$，又向量 $\vec{c}$ 在 $\vec{a},\vec{b}$ 夹角的平分线上，且 $|\vec{c}|=3\sqrt{42}$，求向量 $\vec{c}$.
}

\oneproblem{
    设有一向量与 $x$ 轴正向、$y$ 轴正向的夹角相等，而与 $z$ 轴正向的夹角是前者的两倍，求与该向量同方向的单位向量.
}

\oneproblem{
    设向量 $\vec{a}=\vec{i}+2\vec{j}-\vec{k},\ \vec{b}=-\vec{i}+\vec{j}$.
    \begin{enumerate}[label=(\arabic*)]
        \item 计算 $\vec{a}\cdot\vec{b}$ 及 $\vec{a}\times\vec{b}$；
        \item 求它们夹角 $\theta$ 的正弦与余弦；
        \item 求垂直两向量所在平面的单位向量.
    \end{enumerate}
}

\oneproblem{
    已知一四面体的顶点为 $A_k(x_k,y_k,z_k)\ (k=1,2,3,4)$，求该四面体的体积.
}

\oneproblem{
    判定四点 $A(1,1,1),B(4,5,6),C(2,3,3),D(10,15,17)$ 是否共面？
}

\oneproblem{
    设向量 $\mathbf{a},\mathbf{b},\mathbf{c}$ 不共面，且 $\mathbf{d}=\alpha\mathbf{a}+\beta\mathbf{b}+\gamma\mathbf{c}$，如果 $\mathbf{a},\mathbf{b},\mathbf{c},\mathbf{d}$ 有公共起点.
    \begin{enumerate}[label=(\arabic*)]
        \item 问系数 $\alpha,\beta,\gamma$ 应满足什么条件，才能使向量 $\mathbf{a},\mathbf{b},\mathbf{c},\mathbf{d}$ 的终点在同一平面上？
        \item 如果 $\mathbf{a}=(1,2,1),\ \mathbf{b}=(0,3,1),\ \mathbf{c}=(2,0,3)$，判定向量 $\mathbf{a},\mathbf{b},\mathbf{c}$ 是否共面？
        \item 设 $\mathbf{d}=(-1,-3,1)$，由(2)求 $\alpha,\beta,\gamma$，使得 $\mathbf{d}=\alpha\mathbf{a}+\beta\mathbf{b}+\gamma\mathbf{c}$，如果 $\mathbf{a},\mathbf{b},\mathbf{c},\mathbf{d}$ 有公共起点，它们是否共面？
    \end{enumerate}
}

\oneproblem{
    问当 $t$ 为何值时，空间四点 $A(1,0,0),B(0,2,0),C(0,0,3),D(-1,2,t)$ 共面？并求平面四边形 $ABCD$ 的面积.
}

\oneproblem{
    设 $\vec{a},\vec{b}$ 为两个非零向量，$|\vec{b}|=1,\ (\vec{a},\vec{b})=\dfrac{\pi}{3}$，计算极限 $\displaystyle \lim_{x\to0}\dfrac{|\vec{a}+x\vec{b}|-|\vec{a}|}{x}$.
}

\oneproblem{
    已知向量 $\overrightarrow{AB}=\vec{a},\ \overrightarrow{AC}=\vec{b},\ \angle ADB=\dfrac{\pi}{2}$.
    \begin{enumerate}[label=(\arabic*)]
        \item 证明 $\triangle BAD$ 的面积 $S_{\triangle BAD}=\dfrac{|\vec{a}\cdot\vec{b}||\vec{a}\times\vec{b}|}{2|\vec{b}|^2}$.
        \item 当 $\vec{a},\vec{b}$ 间的夹角为何值时，$\triangle BAD$ 的面积最大，并求最大面积值.
    \end{enumerate}
}

\oneproblem{
    证明：$\mathbf{a},\mathbf{b},\mathbf{c}$ 不共面当且仅当 $\mathbf{a}\times\mathbf{b},\ \mathbf{b}\times\mathbf{c},\ \mathbf{c}\times\mathbf{a}$ 不共面.
}

\oneproblem{
    设 $\mathbf{e}_1,\mathbf{e}_2,\mathbf{e}_3$ 不共面，证明：任一向量 $\mathbf{a}$ 可以表示成
    $
    \mathbf{a}=\dfrac{1}{(\mathbf{e}_1,\mathbf{e}_2,\mathbf{e}_3)}\left[(\mathbf{a},\mathbf{e}_2,\mathbf{e}_3)\mathbf{e}_1+(\mathbf{a},\mathbf{e}_3,\mathbf{e}_1)\mathbf{e}_2+(\mathbf{a},\mathbf{e}_1,\mathbf{e}_2)\mathbf{e}_3\right].
    $
}

\oneproblem{
    设 $\mathbf{a},\mathbf{b},\mathbf{c}$ 不共面，且向量 $\mathbf{r}$ 满足
    $
    \mathbf{a}\cdot\mathbf{r}=\alpha,\ \mathbf{b}\cdot\mathbf{r}=\beta,\ \mathbf{c}\cdot\mathbf{r}=\gamma,
    $
    那么有 \[\mathbf{r}=\dfrac{1}{(\mathbf{a},\mathbf{b},\mathbf{c})}\left[\alpha(\mathbf{b}\times\mathbf{c})+\beta(\mathbf{c}\times\mathbf{a})+\gamma(\mathbf{a}\times\mathbf{b})\right]\]
}

\oneproblem{
    证明：$(\mathbf{a}\cdot\mathbf{b})^2+(\mathbf{a}\times\mathbf{b})^2=|\mathbf{a}|^2|\mathbf{b}|^2$。由此推导用三角形三边长 $a,b,c$ 计算三角形的面积公式，其中向量 $(\mathbf{a})^2=\mathbf{a}\cdot\mathbf{a}$.
}

% ======================================================
% 第 49 天
% ======================================================
\setday{49}{空间平面与直线}

\oneproblem{
    设一平面经过原点及点 $(6,-3,2)$，且与平面 $4x-y+2z=8$ 垂直，求此平面的方程.
}

\oneproblem{
    求与原点的距离为 6，且在三个坐标轴上的截距之比为 $a:b:c=1:3:2$ 的平面方程.
}

\oneproblem{
    已知两个平面 $\pi_1:x+2y-z+1=0$ 和 $\pi_2:2x-y+2z-1=0$，求：
    \begin{enumerate}[label=(\arabic*)]
        \item 这两个平面的夹角 $\theta$ 的余弦；
        \item 这两个平面的角平分面的方程.
    \end{enumerate}
}

\oneproblem{
    已知直线 $L_1$ 和 $L_2$ 的方程 $L_1:\dfrac{x-1}{1}=\dfrac{y-2}{0}=\dfrac{z-3}{-1}$ 和 $L_2:\dfrac{x+2}{2}=\dfrac{y-1}{1}=\dfrac{z}{1}$，试求过 $L_1$ 且平行于 $L_2$ 的平面方程.
}

\oneproblem{
    求过点 $M_0(1,0,1)$ 且与直线 $L:x-1=y+1=z-1$ 垂直相交的直线方程.
}

\oneproblem{
    已知直线 $L:\begin{cases} 2x-4y+z=0 \\ 3x-y-2z=9 \end{cases}$ 和平面 $\pi:4x-y+z=1$，试求直线 $L$ 在平面 $\pi$ 上的投影直线方程.
}

\oneproblem{
    求通过直线 $L:\begin{cases} 2x+y-3z+2=0 \\ 5x+5y-4z+3=0 \end{cases}$ 的两个相互垂直的平面 $\pi_1,\pi_2$，使其中一个平面过点 $(4,-3,1)$.
}

\oneproblem{
    设有空间中五点 $A(1,0,1),B(1,1,2),C(1,-1,-2),D(3,1,0),E(3,1,2)$。试求过点 $E$ 且与 $A,B,C$ 所在平面 $\Sigma$ 平行而与直线 $AD$ 垂直的直线方程.
}

\oneproblem{
    设平面 $\pi$ 方程为 $Ax+By+Cz+D=0$，则向量 $\vec{r}=(r_1,r_2,r_3)$ 平行于平面 $\pi$ 或在平面 $\pi$ 上的充分必要条件是 $Ar_1+Br_2+Cr_3=0$.
}

\oneproblem{
    设相交于直线 $L$ 的两个平面 $\pi_1$ 和 $\pi_2$ 的方程分别为 $A_1x+B_1y+C_1z+D_1=0,\ A_2x+B_2y+C_2z+D_2=0$，则平面 $\pi$ 过直线 $L$（有轴平面束）当且仅当 $\pi$ 的方程形如 $\lambda(A_1x+B_1y+C_1z+D_1)+\mu(A_2x+B_2y+C_2z+D_2)=0$ (*) 其中 $\lambda,\mu$ 是不全为 0 的实数.
}

% ======================================================
% 第 50 天
% ======================================================
\setday{50}{空间平面、直线的方程及位置关系}

\oneproblem{
    试判定直线 $L:\begin{cases} x+3y+2z+1=0 \\ 2x-y-10z+3=0 \end{cases}$ 与平面 $\pi:4x-2y+z-2=0$ 的位置关系.
}

\oneproblem{
    试计算点 $(2,1,0)$ 到平面 $3x+4y+5z=0$ 的距离.
}

\oneproblem{
    求直线 $l_1:\begin{cases} x-y=0 \\ z=0 \end{cases}$ 与直线 $l_2:\dfrac{x-2}{4}=\dfrac{y-1}{-2}=\dfrac{z-3}{-1}$ 的距离.
}

\oneproblem{
    设点 $P(3,1,-4)$ 是直线 $L:\dfrac{x+1}{2}=\dfrac{y-4}{-2}=\dfrac{z-1}{1}$ 外一点。求点 $P$ 在直线 $L$ 上垂足 $Q$ 的坐标，并求点 $P$ 到直线 $L$ 的距离.
}

\oneproblem{
    已知直线 $L$ 通过一平面 $x+y-z-8=0$ 与直线 $\dfrac{x-1}{2}=\dfrac{y-2}{-1}=\dfrac{z+1}{2}$ 的交点，且与直线 $\dfrac{x}{2}=\dfrac{y-1}{1}=\dfrac{z}{1}$ 垂直相交，求直线 $L$ 的方程.
}

\oneproblem{
    求两直线的公垂线的方程.
    $
    L_1:\begin{cases} x+2y+5=0 \\ 2y-z-4=0 \end{cases},\quad L_2:\begin{cases} y=0 \\ x+2z+4=0 \end{cases}
    $
}

\oneproblem{
    设直线 $L$ 过点 $P(-3,5,-9)$ 且与直线 $L_1:\begin{cases} y=4x-7 \\ z=5x+10 \end{cases}$ 与 $L_2:\begin{cases} y=3x+5 \\ z=2x-3 \end{cases}$ 都相交，求直线 $L$ 的方程.
}

\oneproblem{
    设直线 $L$ 过点 $A(1,2,1)$ 且垂直于直线 $L_1:\dfrac{x-1}{3}=\dfrac{y}{2}=\dfrac{z+1}{1}$，又和直线 $L_2:\dfrac{x}{2}=y=\dfrac{z}{-1}$ 相交，求直线 $L$ 的方程.
}

\oneproblem{
    已知空间的两条直线：
    $
    l_1:\dfrac{x-4}{1}=\dfrac{y-3}{-2}=\dfrac{z-8}{1},\quad l_2:\dfrac{x+1}{7}=\dfrac{y+1}{-6}=\dfrac{z+1}{1},
    $
    \begin{enumerate}[label=(\arabic*)]
        \item 证明 $l_1$ 和 $l_2$ 异面；
        \item 求 $l_1$ 和 $l_2$ 公垂线的标准方程；
        \item 求连接 $l_1$ 上任一点和 $l_2$ 上的任一点线段中点的轨迹的一般方程.
    \end{enumerate}
}

\oneproblem{
    求通过直线 $L:\begin{cases} 2x+y-3z+2=0 \\ 5x+5y-4z+3=0 \end{cases}$ 的两个相互垂直的平面 $\pi_1,\pi_2$，使其中一个平面过点 $(4,-3,1)$.
}

\oneproblem{
    设有空间中五点 $A(1,0,1),B(1,1,2),C(1,-1,-2),D(3,1,0),E(3,1,2)$。试求过点 $E$ 且与 $A,B,C$ 所在平面 $\Sigma$ 平行而与直线 $AD$ 垂直的直线方程.
}

\oneproblem{
    证明两直线 $L_1:\begin{cases} A_1x+B_1y+C_1z+D_1=0 \\ A_2x+B_2y+C_2z+D_2=0 \end{cases}$ 和 $L_2:\begin{cases} A_3x+B_3y+C_3z+D_3=0 \\ A_4x+B_4y+C_4z+D_4=0 \end{cases}$ 在同一平面上的充要条件是
    $
    \Delta=\begin{vmatrix}
    A_1 & B_1 & C_1 & D_1 \\
    A_2 & B_2 & C_2 & D_2 \\
    A_3 & B_3 & C_3 & D_3 \\
    A_4 & B_4 & C_4 & D_4
    \end{vmatrix}=0.
    $
}

\oneproblem{
    设 $L_1$ 和 $L_2$ 是空间中两异面直线。设在标准直角坐标系下直线 $L_1$ 过坐标为 $a$ 的点，以单位向量 $v$ 为直线方向；直线 $L_2$ 过坐标为 $b$ 的点，以单位向量 $w$ 为直线方向；
    \begin{enumerate}[label=(\arabic*)]
        \item 证明：存在唯一点 $P\in L_1$ 和 $Q\in L_2$ 使得两点连线 $PQ$ 同时垂直于 $L_1$ 和 $L_2$.
        \item 求 $P$ 点和 $Q$ 点的坐标（用 $a,b,v,w$ 表示）.
    \end{enumerate}
}